% !PRINTED 96
\weberpage{96}

\begin{center}
\textbf{\S\ 28.\quad\gesperrt{Zusammensetzung der Transformationen.}}
\end{center}

Eine Transformation [\S\ 27, (9)] ist vollständig bestimmt durch die Transformationszahlen $a, b, c, \partial$, und diese vier ganzen Zahlen können beliebig gegeben sein, wenn nur ihre Determinante $n = a\partial - bc$ \gesperrt{positiv} ist. Gibt man diesen vier Zahlen das entgegengesetzte Zeichen, so gehen $\omega_1', \omega_2'$ nach \S\ 27, (6) in $-\omega_1', -\omega_2'$ über, und ersetzt man $a, b, c, d$ durch $ma, mb, mc, m\partial$, % [sic] \textquotedbl{}d\textquotedbl{} here, \textquotedbl{}\partial\textquotedbl{} a few words later — same variable printed two ways
worin $m$ eine beliebige natürliche Zahl ist, so gehen $\omega_1', \omega_2'$ in $\omega_1'/m$, $\omega_2'/m$ und $n$ in $m^2n$ über. Das Periodenverhältnis $\omega' = \omega_2'/\omega_1'$ bleibt in diesen beiden Fällen ungeändert. Einstweilen wollen wir aber zwei Transformationen immer als verschieden betrachten, wenn die Transformationszahlen verschieden sind. Nach dieser Festsetzung können wir eine Transformation unzweideutig durch eine Matrix

\weq{1}{S = \begin{pmatrix} a, & b \\ c, & \partial \end{pmatrix}}

darstellen. Die Determinante

\weq{2}{n = a\partial - bc}

ist der Transformationsgrad.

Nach dieser Bezeichnung stellen wir die Relationen (6), \S\ 27 auch so dar:

\weq{3}{n(\omega_1', \omega_2') = \begin{pmatrix} a, & b \\ c, & \partial \end{pmatrix}(\omega_1, \omega_2).}

Setzt man
\[
S' = \begin{pmatrix} a', & b' \\ c', & \partial' \end{pmatrix}, \quad a'\partial' - b'c' = n',
\]
so ist

\weq{4}{n'(\omega_1'', \omega_2'') = \begin{pmatrix} a', & b' \\ c', & \partial' \end{pmatrix}(\omega_1', \omega_2'),}

und wenn man in (4) $\omega_1', \omega_2'$ nach (3) durch $\omega_1, \omega_2$ ausdrückt, so erhält man

\weq{5}{nn'(\omega_1'', \omega_2'') = \begin{pmatrix} a'a+b'c, & a'b+b'\partial \\ c'a+\partial'c, & c'b+\partial'\partial \end{pmatrix}(\omega_1, \omega_2).}

Setzen wir also

\weq{6}{S'' = S'S,}

so ist
\[
S'' = \begin{pmatrix} a'a+b'c, & a'b+b'\partial \\ c'a+\partial'c, & c'b+\partial'\partial \end{pmatrix} = \begin{pmatrix} a'', & b'' \\ c'', & \partial'' \end{pmatrix}, \quad a''\partial'' - b''c'' = n'' = nn'.
\]
